\section{Main findings}
\label{sec:intro}

\subsection{The problem}

A workspace is a repository holding many independently versioned packages that
depend on one another. Releasing it means choosing a new version for each
package that needs one, writing those versions into manifests, uploading
artefacts to registries, and recording what happened. The recording is the
hard part: registries have no transactions, a publish is an irreversible write
to a service the release tool does not control, several registries may be
involved, and no coordinator can commit or abort them as a unit. A run that
publishes seven packages is seven independent, individually fallible writes.

CCME\footnote{\emph{Conventional Commits: Monorepo Extension}, version
3.1.0-rc.2 at commit \texttt{252216ba}, \code{specs/ccme-spec/SPEC.md} in the
\code{dispat} repository. Referred to throughout as \SPEC{}, cited by section
number.} answers this the way a saga does. The release plan is a pure function
of the repository at a fixed commit; each completed leg is recorded by an
annotated git tag written immediately after that leg succeeds; and recovery is
re-execution: run the engine again and it recomputes the same transaction,
skipping the legs already recorded. That design works only if re-execution is
well behaved, and \spec{13.7c} states the eight properties that say it is:

\begin{description}[leftmargin=3.2em,style=nextline,font=\normalfont\bfseries]
  \item[\G{1} termination] propagation halts, whatever the declared depth.
  \item[\G{2} completeness] a contribution admitted to a dependent survives in
    every later run until that dependent releases at or after a release of
    every owed source carrying it; there are no orphans, whether the provider
    failed or the consumer got ahead of it.
  \item[\G{3} version stability] a package released late is released at the
    version it was originally planned at.
  \item[\G{4} no double ledger delivery] a contribution does not survive its
    delivery in the tag ledger; this says nothing unconditional about external
    effects.
  \item[\G{5} no blast-radius widening] under the retry invariant, later
    targets are a subset of the first run's.
  \item[\G{6} restricted convergence] under the retry invariant, repeated
    successful progress empties the initial plan within as many runs as it had
    publishable packages.
  \item[\G{7} channel discharge] a prerelease channel is proposed to a package
    only until it is released on it.
  \item[\G{8} one-way axis independence] suppressing a version bump does not
    suppress a channel change; resolved channels may suppress bumps.
\end{description}

\subsection{What this report establishes}

The report writes the algorithm down once, formally
(\cref{sec:model,sec:propagation,sec:protocol}), proves the guarantees inside
that model (\cref{sec:guarantees}), and reports what the reference
implementation was observed to do in the campaigns that were executed
(\cref{sec:evaluation}). The findings, each with the result that carries it:

\begin{enumerate}[leftmargin=2em]
  \item \textbf{The eight guarantees hold for the model under the explicit
    retry invariant of \spec{13.7c} and two hypotheses the specification does
    not state}, (R) that owed sources reach their target and (K) that
    propagation edges order and block (\cref{thm:g1,thm:g2,thm:g3,thm:g4,thm:g5,thm:g6,thm:g7,thm:g8}).
    (R) holds outright for a unit with a single source; it can fail for a
    multi-source unit, and the reference engine exhibited exactly that failure
    until a fix that postdates the measured builds (\cref{ob:attribution}).
  \item \textbf{Delivery closes the orphan a consumer reaches by getting ahead
    of a failed provider.} The candidate admits a propagated contribution
    until its source has delivered it (\spec{13.4a}); the model transcribes
    this test (\cref{def:delivery}) and \cref{thm:g2} is proved for it. Two
    rules follow and are stated as boundaries rather than proved away: an owed
    window per dependency edge and the refusal \diag{E201}, which the
    specification states and the engine does not yet implement
    (\cref{ob:owed}), and the same-commit case that the protocol reports
    rather than resolves (\cref{ob:samecommit}).
  \item \textbf{The plan is a deterministic, idempotent function of the fixed
    inputs, and a run after a partial failure plans the failed packages, the
    blocked ones and the consumers still owed} (\cref{thm:determinism},
    \cref{cor:idempotency}, \cref{thm:resumption}).
  \item \textbf{The staleness audit of \spec{13.7b} never over-reports and
    under-reports in three shapes}, all with a multi-source unit; the second
    shape is the engine defect above (\cref{thm:duality}, \cref{rem:owed-reach}).
  \item \textbf{Build readiness buys wall time in a form that can be stated.}
    With unlimited slots a run lasts its longest path and a weaker relation
    never lengthens it; the chain, fan-out and fan-in formulas are exact; under
    slot budgets a greedy schedule is within the list-scheduling bound and can
    show Graham's anomaly (\cref{sec:exec-durations}). The placement invariance
    of the distributed profile, that worker count, placement and completion
    order do not change the plan, is stated and proved in the extended edition
    (\cref{prop:exec-placement}).
  \item \textbf{The implementation departs from the text in eight recorded
    places}, none of which loses a release silently except the stranded
    consumer of \cref{ob:owed}; \cref{sec:departures} tabulates them with
    their status.
  \item \textbf{The measurements.} In every state of a finite projection and in
    every completed episode, the consumer-relative window retained exactly the
    residual obligations, and a source-relative window would have lost
    obligations in 84 of the 120 designed fault episodes and in 260 of 772
    states (\cref{sec:eval:rq1}). Recorded progress survived repeated failure
    without version drift in 130 of 130 rows, and the restricted oracle agreed
    on all 300 histories of each companion campaign (\cref{sec:eval:rq2}). A
    configured barrier distinguished local build availability from required
    publication, and both a failed and an absent required test blocked release
    (\cref{sec:eval:rq3}). Planning cost rose from 0.017~s at 100 packages to
    about 19~s at 10,000 on a capped host (\cref{sec:eval:rq4}). On a kernel
    build farm, two groups met their public reference times with a second node,
    twice each after one rejected candidate, where the same node alone did
    not; on an LLVM matrix, a release through two workers took 27:48.94
    against 41:12.02 on one node, once per side; a one-heavy-package fixture
    showed the predicted negative case (\cref{sec:eval:rq5}). No speedup of the
    profile or of the engine is claimed beyond what these pairs measured.
\end{enumerate}

\subsection{Scope, and how to read this edition}

The subject is the release computation of \spec{13} and the publish protocol
of \spec{19}, with the execution rules of \specs{19.2a}{28} as transcribed
obligations. Message syntax (\spec{4}--\spec{8}) is taken as given: units
arrive already parsed. Outside the theorems, and not spoken for by them, are
manifest range reconciliation (\spec{9.4}, \spec{19.5}), the diagnostic
registry of \spec{16} beyond the codes used as labels, executor scheduling
beyond the stated relations and bounds, shared-version groups
(\spec{13.9a}), and the adapter and rollback protocols of \specs{25}{26}. The
report claims no novelty for sagas, idempotent consumers, breadth-first
search or list scheduling; its contribution is the formal model, corrected
hypotheses, counterexamples, the ledger of departures, and the record of what
was measured.

This is the condensed edition. The extended edition, distributed beside it as
\code{technical-report-extended.pdf} and built from the same sources, adds the
transcriptions of the distributed profile (\spec{28}) and of the
repository-local windows (\spec{27}), the cost analysis of
\specs{13.11}{20.7}, the historical audit of CCME 1.0.0 with its six
counterexamples, the conformance table that maps every result to the
specification's vectors and to the integration test behind it, and the full
narratives of the campaigns, including every discarded attempt. A reference
whose number carries an F, such as \emph{Proposition F73}, points into the
extended edition; every other reference is to this document.
\Cref{appendix:notation} is the symbol table, and \cref{appendix:abstract}
reads the companion conference abstract against this report sentence by
sentence. \SPEC{} is licensed GPL-3.0-or-later and this report reproduces its
definitions and pseudocode; it is distributed under the same licence
(\code{paper/LICENSE}).
